\documentclass[11pt]{article}
\usepackage[T1]{fontenc}
\usepackage{lmodern,amsmath,amssymb,amsthm,mathtools,booktabs,microtype}
\usepackage[margin=1.02in]{geometry}
\usepackage{xcolor}
\usepackage[colorlinks=true,linkcolor=blue!45!black,urlcolor=blue!45!black,citecolor=blue!45!black]{hyperref}
\hypersetup{pdftitle={Rigidity at the moment barrier in even-step abelian graph towers},pdfsubject={Spanning-tree valuations and finite-field moment rigidity}}
\setlength{\emergencystretch}{2em}
\newtheorem{theorem}{Theorem}[section]
\newtheorem{lemma}[theorem]{Lemma}
\newtheorem{corollary}[theorem]{Corollary}
\newtheorem{proposition}[theorem]{Proposition}
\theoremstyle{remark}\newtheorem{remark}[theorem]{Remark}
\newcommand{\F}{\mathbb F}
\newcommand{\Z}{\mathbb Z}
\newcommand{\Pj}{\mathbb P}
\newcommand{\A}{\mathcal A}
\newcommand{\ord}{\operatorname{ord}}
\newcommand{\Li}{\operatorname{Li}}
\newcommand{\vp}{v_p}
\DeclareMathOperator{\ev}{ev}
\title{\textbf{Rigidity at the moment barrier in even-step\\ abelian graph towers}\\[5pt]\large Exact spanning-tree extrema in a new congruence class\\ and an asymptotic classification at every large prime}
\author{Research extension of the ProveIt GRH program}
\date{27 November 2026}
\begin{document}
\maketitle
\begin{abstract}
For a rank-two $p$-adic Cayley tower with $m=2s\ge6$ steps and $p>2m$, write $M_j(U)=\sum_i(U\cdot v_i)^{2j}$ for the directional even moments. We prove that $M_1=\cdots=M_{s-1}=0$ holds if and only if the steps consist of two independent radial blocks, each of length $s$, with squared coefficients forming a coset of the $s$th roots of unity. In particular the maximal first nonzero moment index $s$ is attainable precisely when $m\mid p-1$, and no spanning list of rank at least three can attain it. The proof uses homogeneous evaluation spaces, self-duality at a rank barrier, and a Vandermonde argument. Combining this rigidity with the exact graph moment formula resolves an infinite family of cases of the extremal-moment problem posed in the ProveIt article: if $p\equiv 1+m\pmod{2m}$ and $2(p+1)>m(m-2)$, the largest spanning-tree coefficient is exactly $m(p+1)$, rather than the previously available upper bound $m(p+1)+m(m-2)$. All maximizing residue lists are classified in this regime. For $m\nmid p-1$, a smooth complete-intersection curve gives simple graphs attaining the sharp leading slope $m-2$ for all sufficiently large $p$. Thus the leading slope is determined for every large prime. We also give an exact count of maximal-depth row planes and finite arithmetic checks. All graph and rigidity results are unconditional; Dirichlet GRH yields a quantitative prime-counting corollary.
\end{abstract}

\section{Context, scope and main result}
The ProveIt project \cite{project,article,graph,source} consolidates a large collection of transfers from the Dirichlet GRH development \cite{upstream}. Its source section describes a recurring method: replace an analytic or arithmetic constraint by a family of finite-dimensional conditions, then use rank, spanning and determinant constraints to force rigidity. The graph section has an exact reduction of spanning-tree valuations to finite-field moments (\texttt{gra:thm:moments}), an interpolation bound (\texttt{gra:thm:extremal}), and an explicit open question on the maximum for even $m\ge6$ when $2m\nmid p-1$ (\texttt{gra:sec:open}, item 5).

Here we develop the rank part of that method. We do not assume any nonabelian GRH, nor do we need Dirichlet GRH for the new rigidity and graph theorems. The only analytic use of Dirichlet GRH is marked explicitly in Section~\ref{sec:dgrh}. The vendored Lean development reports a proof of the Dirichlet statements at revision \texttt{72fead102f663d5c7a407964a13dc0cf9d977207}; this report does not rerun its build or certify its proof. The merged article's README says it is a merge, not an independent verification. We use its graph question as motivation and rederive the particular graph formula needed here. The research direction also relates to the Iwasawa theory of graph towers, notably \cite{dubose,dion}.

Let $p$ be a prime, $m=2s\ge6$, and $p>2m$. A list $S=(v_1,\dots,v_m)$ in $\F_p^2$ must span $\F_p^2$; zeros and repetitions are permitted unless we specify a simple list (nonzero and distinct up to sign). For a nonzero covector $u$, let $r_S([u])$ be the first $j\ge1$ with $M_j(u)\ne0$, and set
\[
 \A(S)=2\sum_{L\in\Pj^1(\F_p)} r_S(L),\qquad
 \A_{\max}(p,m)=\max_S\A(S).
\]
Let $t(S)$ be the first index with $M_j$ a nonzero polynomial. The maximum over simple lists is denoted by $\A_{\max}^{\rm simp}$.

\begin{theorem}[Rigidity and extremal laws]\label{thm:main}
Under these hypotheses:
\begin{enumerate}
\item $t(S)\le s$. Equality holds exactly for lists of the form
\begin{equation}\label{eq:block}
 (\alpha_1 e,\ldots,\alpha_s e,\beta_1 f,\ldots,\beta_s f),
\quad \{\alpha_i^2\}=a^2\mu_s,\quad \{\beta_i^2\}=b^2\mu_s,
\end{equation}
where $e,f$ are independent, $a,b\ne0$, and $\mu_s\subset\F_p^\times$ is the group of $s$th roots of unity of cardinality $s$. This is possible if and only if $m\mid p-1$.
\item If $2m\mid p-1$, then $\A_{\max}=\A_{\max}^{\rm simp}=m(p+1+m)$ (the numerical bound was already known), and \eqref{eq:block} classifies all maximizers.
\item If $m\mid p-1$ but $2m\nmid p-1$, then
\begin{equation}\label{eq:middle}
 m(p+1)\le\A_{\max}\le\max\{m(p+1),(m-2)(p+1+m)\}.
\end{equation}
If $2(p+1)>m(m-2)$, both maxima equal $m(p+1)$, and \eqref{eq:block} classifies all maximizers.
\item If $m\nmid p-1$, then
\begin{equation}\label{eq:low}
 \A_{\max}\le(m-2)(p+1+m).
\end{equation}
For all sufficiently large such $p$ (with explicit cutoff in Theorem~\ref{thm:curve}),
\begin{equation}\label{eq:asy}
 (m-2)(p+1)\le\A_{\max}^{\rm simp}\le\A_{\max}\le(m-2)(p+1)+m(m-2).
\end{equation}
\end{enumerate}
\end{theorem}
The exact law in part 3 and the slope drop in part 4 are the new extremal consequences. The gap in part 4 is bounded independently of $p$ but is not generally determined here. The statement about all maximizers is about residue lists; integral lifts can vary freely. The classification is stronger than an example attaining a bound, and requires no finite search.

\section{The finite-field rigidity theorem}\label{sec:rigidity}
Work first over any field $K$ of characteristic zero or of characteristic greater than $4s$. Let $v_i=(x_i,y_i)\in K^2$ span $K^2$. Let $R$ be the number of distinct radial lines among the nonzero steps, and let $n_\ell$ count the nonzero steps on line $\ell$. In particular $R\ge2$ and $\sum_\ell n_\ell\le2s$. Assume for the rest of this section that $M_1,\ldots,M_{s-1}$ vanish identically. Invertibility of binomial coefficients gives the mixed-moment identities
\begin{equation}\label{eq:mixed}
 \sum_{i=1}^{2s} H(v_i)=0
 \quad\text{for every homogeneous }H\in K[X,Y]_{2j},\quad 1\le j<s.
\end{equation}
For $k\ge1$ define the evaluation space
\[
 E_k=\{(H(v_1),\dots,H(v_{2s})):H\in K[X,Y]_k\}\subseteq K^{2s}.
\]
Use the nondegenerate symmetric bilinear form $\langle a,b\rangle=\sum_i a_i b_i$; this is not a Hermitian form.

\begin{lemma}[Evaluation rank and orthogonality]\label{lem:rank}
For $k\ge1$, $\dim E_k=\min(R,k+1)$. If $k+l$ is even and $2\le k+l\le2s-2$, then $E_k\perp E_l$. In particular $E_{s-1}$ is totally isotropic. When $s$ is odd, $E_{s-1}\perp\mathbf1$ for $\mathbf1=(1,\ldots,1)$.
\end{lemma}
\begin{proof}
Choose one representative $w_\ell$ per radial line. At $v_i=c_iw_\ell$ ($c_i\ne0$) the evaluation is $c_i^k H(w_\ell)$; zero columns contribute zero. The evaluation rank therefore equals the rank at $R$ distinct projective points. A nonzero homogeneous form of degree $k$ in two variables has at most $k$ distinct projective roots; moreover when $R\le k+1$, products of linear forms vanishing at all but one of the $R$ points give interpolants (multiply by further linear factors nonvanishing at that point if necessary). This proves the rank formula. Pairings are sums of evaluations of products, so \eqref{eq:mixed} proves the orthogonality. The final assertion uses $s-1$ even and positive.
\end{proof}

\begin{lemma}[Too many radial lines are impossible]\label{lem:many}
If $s$ is odd, then $R<s$. If $s$ is even (hence $s\ge4$), then $R\le s$.
\end{lemma}
\begin{proof}
If $R\ge s$, then $\dim E_{s-1}=s$. Total isotropy in the nondegenerate space of dimension $2s$ gives $E_{s-1}=E_{s-1}^\perp$.

For odd $s$, Lemma~\ref{lem:rank} implies $\mathbf1\in E_{s-1}$. But $\langle\mathbf1,\mathbf1\rangle=2s\ne0$, contradicting isotropy.

For even $s\ge4$, suppose $R>s$. By the same lemma $E_1\perp E_{s-1}$, since $1+s-1=s$ is even and $s\le2s-2$. Hence $E_1\subseteq E_{s-1}$. There are homogeneous forms $P,Q$ of degree $s-1$ with $P(v_i)=x_i$, $Q(v_i)=y_i$ for every $i$. The homogeneous degree-$s$ form $YP-XQ$ vanishes on $R>s$ distinct projective points, so it is identically zero. Since $X,Y$ are coprime, $P=XH$, $Q=YH$ for one homogeneous form $H$ of degree $s-2$. At each nonzero step $H(v_i)=1$, whereas at a zero step $H(0)=0$. The degree $s-2$ is positive and even; \eqref{eq:mixed} gives
\[
 0=\sum_i H(v_i)=\#\{i:v_i\ne0\},
\]
contradicting the nonzero integer count, which is at most $2s$ and hence nonzero in $K$.
\end{proof}

\begin{lemma}[Too few radial lines, except two, are impossible]\label{lem:few}
If $2\le R\le s$, then $R=2$, $n_1=n_2=s$, and in each radial line the first $s-1$ power sums of the squared coefficients vanish.
\end{lemma}
\begin{proof}
Write the nonzero steps on line $\ell$ as $c_{\ell i} w_\ell$. Put $k_0=\lfloor R/2\rfloor\ge1$. For $k_0\le j\le s-1$, the $R$ forms $(u\cdot w_\ell)^{2j}$ are independent: the evaluation matrix of degree-$2j$ at the $w_\ell$ has rank $R\le2j+1$, and multiplication by the invertible binomial coefficients and transposition gives the claimed independence. Since
\[
 0=M_j(u)=\sum_{\ell=1}^R (u\cdot w_\ell)^{2j}\sum_i c_{\ell i}^{2j},
\]
we get $\sum_i c_{\ell i}^{2j}=0$ for every $j=k_0,\ldots,s-1$ and every $\ell$.

If a fixed line has $d$ distinct nonzero squared coefficients $z_1,\ldots,z_d$ with multiplicities $h_1,\ldots,h_d$, then $d$ consecutive zero power sums at positive indices would give a Vandermonde matrix on the $z_a$ times the vector $(h_a z_a^{k_0})_a$. This matrix is invertible, and each entry of the vector is nonzero (the multiplicities are positive and at most $2s$). Thus $d>s-k_0$, and
\begin{equation}\label{eq:radialcount}
 n_\ell\ge s-\lfloor R/2\rfloor+1.
\end{equation}
If $R=3$, this gives at least $3s>2s$ steps. If $4\le R\le s$, it gives at least $4(s-\lfloor s/2\rfloor+1)>2s$. Both are impossible. For $R=2$, each count is at least $s$; hence each is exactly $s$ and no zero steps remain. The vanishing sums for $1\le j<s$ were just proved.
\end{proof}

\begin{theorem}[Classification at the barrier, over general fields]\label{thm:rigidity}
For $m=2s\ge6$, spanning rank-two lists over $K$ satisfying \eqref{eq:mixed} are exactly the lists \eqref{eq:block}, interpreted over $K$. Their first nonzero moment is $M_s$. In particular there are no spanning lists with $M_1=\cdots=M_s=0$.
\end{theorem}
\begin{proof}
Lemmas~\ref{lem:many} and \ref{lem:few} force precisely two blocks of length $s$. If $z_1,\ldots,z_s$ are the squared coefficients in either block, Newton's identities and the first $s-1$ zero power sums yield
\[
 \prod_{i=1}^s(T-z_i)=T^s-c,
 \qquad c\ne0.
\]
Thus the $z_i$ are distinct and are all $z_1$ times the $s$th roots of unity in $K$. Since each $z_i$ is a square, $z_1=a^2$ for some $a$, and the entire root-of-unity group consists of squares. Conversely these two blocks have all required moments zero by the geometric sum formula. The next moment is
\[
 M_s(u)=s\big(a^{2s}(u\cdot e)^{2s}+b^{2s}(u\cdot f)^{2s}\big),
\]
which is a nonzero polynomial because $e,f$ are independent and $s\ne0$ in $K$. This proves all assertions.
\end{proof}

\begin{corollary}[Higher-rank obstruction]\label{cor:rank}
For any spanning list of $2s\ge6$ vectors in $K^d$ with $d\ge3$, at least one of the moments $M_1,\ldots,M_{s-1}$ is nonzero.
\end{corollary}
\begin{proof}
Extend scalars to an algebraic closure. If all the moments vanished, every linear projection to dimension two would preserve their vanishing. Choose a projection nonzero on all nonzero steps and taking every pair of nonproportional steps to independent vectors. Such maps form a nonempty Zariski-open subset of the space of $2\times d$ matrices: for each pair the determinant polynomial is nonzero, and for each nonzero step one may impose a nonzero first-coordinate linear polynomial. The algebraic closure is infinite, so finitely many such polynomials are simultaneously nonzero somewhere. The image spans rank two and preserves the number of radial lines, at least $d\ge3$. Theorem~\ref{thm:rigidity} contradicts this.
\end{proof}

\begin{remark}[Why four steps are exceptional]
For $s=2$, the argument with $H$ in Lemma~\ref{lem:many} would use a constant form, whose sum is not constrained to vanish. The article's four-step isotropic-plane families are therefore not contradicted. The hypotheses $s\ge3$ and characteristic greater than $4s$ matter in the proof.
\end{remark}

\section{Spanning trees and exact extremal coefficients}\label{sec:graph}
We give the needed reduction from the article, to make the result checkable independently of that text. Let integral lifts of the steps be $\widetilde v_i\in\Z^d$, spanning after reduction modulo $p$. Let $G_n$ be the undirected Cayley multigraph on $(\Z/p^n\Z)^d$ with one edge in the positive step direction for each vertex and each listed step (reverse incidence supplies the negative direction). Repeated steps are counted with multiplicity; zero steps contribute no Laplacian term. The group is generated by the steps since they span modulo $p$. Write $\kappa_n$ for the number of spanning trees.

\begin{lemma}[Exact moment formula, all levels]\label{lem:tree}
For $p>2m$ and rank $d$, the depth $r_S(L)$ of every residual projective direction is at most $m$, and
\begin{equation}\label{eq:tree}
 \vp(\kappa_n)=\left(2\sum_{L\in\Pj^{d-1}(\F_p)}r_S(L)\right)
 \sum_{h=0}^{n-1}p^{(d-1)h}-dn.
\end{equation}
\end{lemma}
\begin{proof}
For any nonzero residual covector, its squared projections $z_i$ are not all zero. If the first $m$ power sums vanish, Newton's identities give $\prod_i(T-z_i)=T^m$, a contradiction. Hence the stated depth exists.

At a character of order $p^h$, let $\zeta$ be a primitive $p^h$th root and $a$ a primitive vector modulo $p^h$ giving the character. Its Laplacian eigenvalue is
\[
 F_a(\zeta-1),\qquad F_a(T)=\sum_{i=1}^m\big(2-(1+T)^{a\cdot\widetilde v_i}-(1+T)^{-a\cdot\widetilde v_i}\big).
\]
This integral power series converges at $T=\zeta-1$. Over $\F_p$, the truncated change of variables $1+T=\exp Z$ through degree $2m$ is invertible (all factorials are units). It gives
\[
 F_a(\exp Z-1)=-2\sum_{j=1}^m\frac{M_j(\bar a)}{(2j)!}Z^{2j}\pmod{Z^{2m+1}}.
\]
Thus the first unit coefficient of $F_a$ is at degree $2r_S([\bar a])\le2m$. Normalize $\ord_\pi(\pi)=1$ for $\pi=\zeta-1$, so $\ord_\pi(p)=p^{h-1}(p-1)\ge2m$. The constant coefficient of $F_a$ is exactly zero, and lower positive degrees with coefficients divisible by $p$ have valuation strictly above $2r_S([\bar a])$. Higher degrees also have strictly larger valuation. Even in the boundary case $p-1=2m$, this argument is strict because the constant term is zero. Hence the eigenvalue has $p$-adic valuation $2r_S([\bar a])/\big(p^{h-1}(p-1)\big)$.

At conductor $p^h$, each residual projective direction has $(p-1)p^{d(h-1)}$ primitive vectors reducing into it. Sum the valuations, then apply the character diagonalization of the Laplacian and the matrix-tree theorem, dividing the nontrivial eigenvalue product by $p^{dn}$. This gives \eqref{eq:tree}.
\end{proof}

For rank two, \eqref{eq:tree} specializes to
\begin{equation}\label{eq:alllevels}
 \vp(\kappa_n)=\A(S)\frac{p^n-1}{p-1}-2n.
\end{equation}
Thus any maximum or inequality for $\A$ is simultaneously valid at every level, and depends only on residues. We now prove the extremal part of Theorem~\ref{thm:main}.

\begin{lemma}[Interpolation and depth ceiling]\label{lem:upper}
If $t$ is the first nonzero moment form and $D$ an upper bound for every direction's depth, then
\[
 \A(S)\le 2t(p+1)+4t(D-t).
\]
One may always take $D=m$, and may take $D=m-1$ if $2m\nmid p-1$.
\end{lemma}
\begin{proof}
A nonzero binary form of degree $2t$ has at most $2t$ distinct projective roots; elsewhere the depth is exactly $t$. At the roots it is at most $D$. The bound follows. A direction of depth $m$ has squared projections whose first $m-1$ power sums vanish, and whose $m$th is nonzero. Newton's identities show that they are precisely the distinct roots of $T^m-c$ with $c\ne0$. Their quotients contain all $m$th roots of unity and are squares; the cyclic square subgroup of $\F_p^\times$ has order $(p-1)/2$. Thus $2m\mid p-1$ is necessary. The universal $D=m$ follows from Lemma~\ref{lem:tree}.
\end{proof}

\begin{proof}[Proof of Theorem~\ref{thm:main}, parts 1--4 except the large-prime lower bound]
The classification and $t\le s$ follow from Theorem~\ref{thm:rigidity}. In $\F_p$, the condition that $s$th roots of unity exist and are squares is equivalent to $2s=m\mid p-1$. In coordinates on covectors scaled by $a,b$ of \eqref{eq:block}, the moments of the classified lists are exactly
\begin{equation}\label{eq:momblocks}
 M_j(X,Y)=\begin{cases}0,&s\nmid j,\\s(X^{2j}+Y^{2j}),&s\mid j.\end{cases}
\end{equation}
Every classified list is simple modulo $p$. The $s$th moment vanishes in exactly $m$ projective directions when $2m\mid p-1$ and none when $m\mid p-1$ but $2m\nmid p-1$: indeed $X^m=-Y^m$ requires $Y\ne0$ and $-1$ to be an $m$th power. In the first case the $m$th moment is nonzero at these roots (the intervening moments vanish), so the depths are $s$ and $m$. In the second case the depth is constantly $s$. Consequently
\begin{equation}\label{eq:topA}
 \A(S)=\begin{cases}m(p+1+m),&2m\mid p-1,\\m(p+1),&m\mid p-1,\ 2m\nmid p-1.\end{cases}
\end{equation}
For $t<s$, the interpolation bound of Lemma~\ref{lem:upper} is strictly increasing in $t$ for $1\le t\le s$ and $D\in\{m-1,m\}$: the successive difference is $2(p+1+2D-4t-2)>0$. With $D=m$, the universal maximum bound at $t=s$ equals the first line of \eqref{eq:topA}, and no $t<s$ can attain it. With $D=m-1$, the bound at $t=s-1$ is $(m-2)(p+1+m)$, proving \eqref{eq:middle} and \eqref{eq:low}. The exact middle-case result holds when $m(p+1)>(m-2)(p+1+m)$, equivalent to the stated threshold. Again all maximizers must be classified lists. The remaining lower bound is supplied in Section~\ref{sec:curve}.
\end{proof}

\begin{corollary}[Extrema of every tower layer]\label{cor:layers}
For $p\equiv1+m\pmod{2m}$, $m\ge6$ even, $p>2m$, and $2(p+1)>m(m-2)$,
\[
 \max_S\vp(\kappa_n)=m(p+1)\frac{p^n-1}{p-1}-2n\qquad(n\ge1).
\]
The maximum holds for simple graphs and one fixed residue list attains it at all layers. The equality configurations are exactly the lifts of \eqref{eq:block}. In particular for $m=6$ the formula holds for every prime $p\equiv7\pmod{12}$ with $p>12$; for $m=8$ it holds for every prime $p\equiv9\pmod{16}$ with $p>16$.
\end{corollary}
\begin{proof}
Use \eqref{eq:alllevels}. For six steps the threshold is automatic; for eight steps the first candidate prime in the specified range is $25$, which is composite, and the next integer in that class is $41$.
\end{proof}

\begin{corollary}[Number and orbit of top-depth planes]\label{cor:count}
For $m=2s\ge6$, the number of rank-two row planes $W\subset\F_p^m$ with $t(W)=s$ is zero if $m\nmid p-1$ and otherwise equals
\begin{equation}\label{eq:count}
 \frac12\binom{2s}{s}\left(2^{s-1}(s-1)!\right)^2.
\end{equation}
They form a single orbit under coordinate permutations and independent coordinate sign changes. In the exact regimes of Theorem~\ref{thm:main}, parts 2--3, this also counts all maximizing planes.
\end{corollary}
\begin{proof}
A plane of the classified form has two generating rows with disjoint supports of size $s$ partitioning the coordinates. The two support blocks are uniquely recovered: all other projective row directions have support of size $2s$. For each unordered partition, choose a line on each support. Normalize the first entry in that block to $1$; the other squared entries are a permutation of $\mu_s\setminus\{1\}$, with two independent square-root choices each. This yields exactly $2^{s-1}(s-1)!$ lines per block and gives \eqref{eq:count}. Coordinate permutations and signs act transitively on these choices, including the partition. No quotient by a further symmetry is needed, because the disjoint-support projective lines are unique.
\end{proof}

\section{Sharp leading slope at the other primes}\label{sec:curve}
We next show that the drop in $t$ when $m\nmid p-1$ is exactly one for all sufficiently large primes, not merely at most $s-1$.

\begin{theorem}[Simple lists at the next moment barrier]\label{thm:curve}
Let $s\ge3$, $m=2s$ and $p>4s$ with $m\nmid p-1$. Put
\[
 D_s=2^{s-2}(s-2)!,\qquad g_s=1+\frac{D_s}{2}\big((s-2)(s-1)-s\big).
\]
If $p+1-2g_s\sqrt p>s^2D_s$, then there exists a simple rank-two list with $t=s-1$. More precisely,
\begin{equation}\label{eq:lowerref}
 \A_{\max}^{\rm simp}\ge (m-2)(p+1)+2N,\qquad
 N=\begin{cases}d,&(p-1)/d\text{ even},\\0,&(p-1)/d\text{ odd},\end{cases}
 \quad d=\gcd(m-2,p-1).
\end{equation}
\end{theorem}
\begin{proof}
Consider the projective variety in $\Pj^{s-1}$ defined by
\begin{equation}\label{eq:curve}
 C_s:\quad \sum_{i=1}^s X_i^{2j}=0\quad(1\le j\le s-2).
\end{equation}
We check the geometric input over an algebraic closure of $\F_p$. It is nonempty, since one can take square roots of the $s$ distinct $s$th roots of unity. At a nonzero point of the affine cone, dependence of the $s-2$ gradient rows would give a nonzero polynomial $B(T)$ of degree at most $s-3$ satisfying $X_i B(X_i^2)=0$ for each $i$ (the factors $2j$ are units). There would be at most $s-3$ distinct nonzero squared coordinates. But the first $s-2$ vanishing power sums, by the Vandermonde argument from Lemma~\ref{lem:few}, force their nonzero multiplicities to vanish, impossible. This also covers $s=3$, when a nonzero constant $B$ cannot vanish at a nonzero squared coordinate. Thus $C_s$ is a smooth complete intersection of dimension one, with defining degrees $2,4,\ldots,2s-4$.

For completeness, a positive-dimensional projective complete intersection is geometrically connected: the Koszul resolution of its structure sheaf and the cohomology of line bundles on $\Pj^{s-1}$ give $H^0(C_s,\mathcal O)=\overline{\F}_p$. In detail the negative twists in homological degrees $1$ through $s-2$ have zero cohomology in degrees $0$ through $s-2$, so the sole contribution to degree zero is $H^0(\Pj^{s-1},\mathcal O)$; no differential can kill it, as the constant section restricts to $1$ on a nonempty variety. Smoothness and connectedness imply geometric irreducibility. B\'ezout gives degree $D_s$; adjunction gives $\omega_{C_s}=\mathcal O_{C_s}((s-2)(s-1)-s)$ and hence genus $g_s$. These are standard complete-intersection formulas (see \cite{hartshorne}).

Discard the $s^2$ hyperplanes $X_i=0$, $X_i=\pm X_j$ for $i<j$. None contains the curve, since the above root-of-unity point belongs to their complement. Each hyperplane meets the curve in at most $D_s$ geometric points. The Hasse--Weil bound \cite{weil} gives at least $p+1-2g_s\sqrt p$ rational points, so the strict inequality of the theorem supplies a rational point in the complement. For a representative $w=(w_1,\ldots,w_s)$, each coordinate is nonzero and distinct up to sign.

Use the simple step list $(w_i,0)$ and $(0,w_i)$ for $1\le i\le s$. Write $P_j=\sum_i w_i^{2j}$. We have $P_1=\cdots=P_{s-2}=0$, and $P_{s-1}\ne0$ because otherwise Newton's identities force the squared coordinates to be a root-of-unity block and $m\mid p-1$. Also $P_s\ne0$: Newton's identity at index $s$ reduces to $P_s=(-1)^{s-1}s\prod_i w_i^2$, since all lower power sums vanish except possibly $P_{s-1}$ whose coefficient $e_1=P_1$ vanishes. Here $s\ge3$ is used. Therefore
\[
 M_{s-1}=P_{s-1}(X^{m-2}+Y^{m-2}),\qquad M_s=P_s(X^m+Y^m).
\]
If $X^{m-2}+Y^{m-2}=0$, the ratio $z=X/Y$ is nonzero; simultaneous vanishing of $X^m+Y^m$ would imply $z^2=1$ and $z^{m-2}=-1$, impossible in odd characteristic. Thus all $N$ roots of the first form have depth $s$, and all other directions have depth $s-1$. The cyclic group $\F_p^\times$ gives the root count in \eqref{eq:lowerref}. This proves the lower bound and $t=s-1$.
\end{proof}

Theorem~\ref{thm:curve} finishes Theorem~\ref{thm:main}. For example $s=3$ gives $D_s=2,g_s=0$, so the cutoff holds when $p+1>18$. For $s=4$, $D_s=8,g_s=9$, and the cutoff is $p+1-18\sqrt p>128$. In all dimensions considered here the cutoff is explicit but not optimized; the construction can work much earlier. Combining with \eqref{eq:alllevels} gives all-level valuation bounds for simple and unrestricted lists.

\section{Distribution over primes and Dirichlet GRH}\label{sec:dgrh}
Let $m\ge6$ be fixed and even. The preceding theorems prove, uniformly as $p\to\infty$ through primes,
\begin{equation}\label{eq:slope}
 \frac{\A_{\max}(p,m)}{p+1}=
 \begin{cases}m+O_m(p^{-1}),&p\equiv1\pmod m,\\m-2+O_m(p^{-1}),&p\not\equiv1\pmod m.\end{cases}
\end{equation}
The same holds for the simple maximum. Hence the empirical law of this normalized statistic over primes converges to a two-point distribution assigning mass $1/\varphi(m)$ to $m$ and $1-1/\varphi(m)$ to $m-2$, by the prime number theorem in arithmetic progressions. Within the high-slope class there are two equally dense subclasses modulo $2m$, with exact offsets $m^2$ and zero from $m(p+1)$ for all sufficiently large primes. This qualitative result is unconditional.

\begin{corollary}[Effective abundance under Dirichlet GRH]
Assume Dirichlet GRH as stated in \cite{upstream}. For $X\ge2m$, define the weighted count of primes admitting a list with $t=s$ by
\[
 H_m(X)=\sum_{\substack{2m<p\le X\\\exists S:\,t(S)=s}}\log p.
\]
Then
\[
 H_m(X)=\frac{X}{\varphi(m)}+O\big(\sqrt X\log^2(mX)+m\log(2m)\big).
\]
The two high-slope classes separately have main terms $X/\varphi(2m)$ with the same shape of error. For the exact formula in the middle class, excluding the additional finitely many primes below its threshold changes the weighted count by $O_m(1)$.
\end{corollary}
\begin{proof}
Theorem~\ref{thm:main} makes the existence event exactly $p\equiv1\pmod m$ above $2m$. The standard Dirichlet-GRH explicit-formula estimate for the prime-weighted progression count is $\theta(X;q,a)=X/\varphi(q)+O(\sqrt X\log^2(qX))$ for $(q,a)=1$ (see \cite{lee}); subtract the primes at most $2m$. The same argument for $q=2m$ applies to the two classes $1$ and $1+m$, both coprime to $2m$. No hypothesis on Dedekind zeta functions or Artin $L$-functions is used.
\end{proof}

\section{Reproducible checks and finite six-step results}
The general proof above is independent of computation. The archive supplies two independent styles of check. \texttt{verify.py} tests the block construction over 212 pairs $(p,s)$, checks a few exact norm products in $\Z[\zeta_p]$, and finds example curve points. \texttt{verify\_six.cpp} performs a complete finite search at ten primes for six steps.

In the finite search, an independent pair of steps is normalized to $(1,0),(0,1)$ by a linear change of coordinates and reordering. Every remaining step is represented by the lexicographically smaller of $v,-v$, including the zero class. The remaining four classes are enumerated as a nondecreasing multiset. A meet-in-the-middle table indexes pairs of classes by their three quadratic monomial sums $(\sum x_i^2,\sum x_i y_i,\sum y_i^2)$. Matching complements to $(-1,0,-1)$ exhausts $M_1\equiv0$; every eligible list is considered in at least one such normalization. Each list is evaluated on all $p+1$ projective directions by direct finite-field power sums. The algorithm also detects simplicity. If $M_1$ is not identically zero, Lemma~\ref{lem:upper} bounds $\A$ by $2(p+1)+4(D-1)$, strictly below the maxima found at every tested prime. Thus the reported maxima are exact across all lists, not just the normalized $M_1=0$ candidates. Counts of normalized candidates are not counts of row planes or orbits.

\begin{center}
\begin{tabular}{rrrrr}
\toprule
$p$ & normalized candidates & $t=3$ & $\A_{\max}$ & $\A_{\max}^{\rm simp}$\\
\midrule
13 & 1,130 & 1 & 120 & 120\\
17 & 4,110 & 0 & 88 & 84\\
19 & 6,380 & 1 & 120 & 120\\
23 & 16,519 & 0 & 104 & 104\\
29 & 56,210 & 0 & 128 & 128\\
31 & 73,324 & 1 & 192 & 192\\
37 & 187,449 & 1 & 264 & 264\\
41 & 311,765 & 0 & 184 & 184\\
43 & 377,047 & 1 & 264 & 264\\
47 & 588,626 & 0 & 200 & 200\\
\bottomrule
\end{tabular}
\end{center}
The search is in particular an exact certificate by exhaustive arithmetic for the five tested primes with $6\nmid p-1$, not a general formula for that class. At $p=19$, the list $\{(1,0),(7,0),(8,0),(0,1),(0,7),(0,8)\}$ gives $\A=120$; the level-one tree valuation is $118$. At $p=29$ a simple list with extra steps $(0,2),(0,12),(2,3),(13,4)$ beyond the independent pair gives $\A=128$. Exact determinants of multiplication matrices in $\Z[Z]/\Phi_p(Z)$ check the level-one valuations $118,118,102,126$ at $p=13,19,23,29$, respectively. This checks the norm product rather than approximating complex eigenvalues.

\section{What remains open and further research questions}
The explicit open extremal question in the graph section of the repository is partly settled: the non-$2m$ congruence class with $m\mid p-1$ now has an exact answer for every sufficiently large prime and a complete equality classification. The first nonzero moment index and leading asymptotic slope are settled for every sufficiently large prime. The remaining bounded offset in \eqref{eq:asy}, and finitely many smaller primes for each fixed $m$, are not resolved by this proof. In particular the work does not claim to resolve a classical conjecture such as the Artin primitive-root conjecture, general automorphic GRH or a Millennium Prize problem. The source article's previously known $2m\mid p-1$ numerical maximum has not been relabeled as new.

\begin{enumerate}
\item Determine the exact bounded offset when $m\nmid p-1$. Even the $m=6$ computations show several offsets and a possible difference between simple and unrestricted graphs. A next step is to classify the two-dimensional spaces annihilating the first $s-2$ moments and study the common roots of their remaining moments.
\item Improve the threshold in the middle class; can the exact answer $m(p+1)$ be established at every $p>2m$ in that class? The interpolation comparison is not strong enough near $p=2m$ for all $m$.
\item Find the largest first nonzero moment index in rank $d\ge3$ for arbitrary $m$. Corollary~\ref{cor:rank} rules out the top rank-two barrier, but does not determine the next barrier in higher rank.
\item Improve the finite-field existence cutoff and the lower bound using two different blocks on $C_s$, rather than identical blocks. The exceptional directions where several moments vanish then become a problem of simultaneous algebraic equations.
\item Formalize the evaluation-rank, isotropic-complement and radial-block lemmas in Lean; the graph transfer needs Newton identities and the cyclotomic/matrix-tree formula. The Hasse--Weil step is separate. No Lean proof or kernel-check claim is supplied with this paper.
\end{enumerate}

\paragraph{Novelty and provenance.} The new contribution is the self-duality and homogeneous interpolation argument, the resulting classification, and its extremal and distribution consequences. We inspected the project README, article outline, source mechanism, the graph section, selected related pending report READMEs, and the vendored theorem statements. The explicit question in the inspected graph section has not already been answered there in the form given above. The literature comparison \cite{dubose,dion} is targeted, not an exhaustive priority search. The new theorems are candidate research contributions pending independent review. Neither the upstream proof nor all 194 merged reports and 15 pending packets were independently audited. We did not modify the repository.

\begin{thebibliography}{9}\small
\bibitem{project} ProveIt, \emph{Generalized Riemann Hypothesis}, repository README, inspected 27 November 2026. \url{https://github.com/VladimirReshetnikov/ProveIt/tree/main/Analysis/GeneralizedRiemannHypothesis}.
\bibitem{article} ProveIt, \emph{Beyond the Riemann hypothesis for Dirichlet L-functions: consequences, transfers and limits}, article README (194 merged reports). \url{https://github.com/VladimirReshetnikov/ProveIt/tree/main/Analysis/GeneralizedRiemannHypothesis/Article}.
\bibitem{graph} ProveIt, article section \emph{Graph towers, spanning-tree valuations and signed cycles}, source \texttt{gra.tex}, blob \texttt{6a866b41da2111c5fc7f9bd10dc4315c9a283de2}. \href{https://github.com/VladimirReshetnikov/ProveIt/blob/main/Analysis/GeneralizedRiemannHypothesis/Article/sections/gra.tex}{Repository source}. In particular labels \texttt{gra:thm:moments}, \texttt{gra:thm:extremal}, \texttt{gra:sec:open}.
\bibitem{source} ProveIt, article section \emph{The source theorem and its transferable mechanism}, source \texttt{src.tex}, blob \texttt{f69668b2f95261630bb1282c2cc1f9c172b7e530}. \href{https://github.com/VladimirReshetnikov/ProveIt/blob/main/Analysis/GeneralizedRiemannHypothesis/Article/sections/src.tex}{Repository source}.
\bibitem{upstream} L. Sernau et al., \emph{The generalized Riemann hypothesis for Dirichlet L-functions, in Lean 4}, revision \texttt{72fead102f663d5c7a407964a13dc0cf9d977207}. \url{https://github.com/lsernau/grh-lean/tree/72fead102f663d5c7a407964a13dc0cf9d977207/grh}. See \texttt{src/Statements.lean}.
\bibitem{dubose} S. DuBose and D. Valli\`eres, On $\Z_\ell^d$-towers of graphs, \emph{Algebraic Combinatorics} \textbf{6} (2023), 1331--1346. \href{https://doi.org/10.5802/alco.304}{doi:10.5802/alco.304}.
\bibitem{dion} C. Dion, A. Lei, A. Ray and D. Valli\`eres, On the distribution of Iwasawa invariants associated to multigraphs, \emph{Nagoya Math. J.} \textbf{253} (2024), 48--90. \href{https://doi.org/10.1017/nmj.2023.18}{doi:10.1017/nmj.2023.18}.
\bibitem{hartshorne} R. Hartshorne, \emph{Algebraic Geometry}, Springer, 1977, Chapters II--III (adjunction, projective-space cohomology and complete intersections).
\bibitem{weil} A. Weil, \emph{Sur les courbes alg\'ebriques et les vari\'et\'es qui s'en d\'eduisent}, Hermann, 1948 (Hasse--Weil bound for curves).
\bibitem{lee} E. S. Lee, The prime number theorem for primes in arithmetic progressions at large values, \emph{Q. J. Math.} \textbf{74} (2023), 1505--1533. \href{https://doi.org/10.1093/qmath/haad031}{doi:10.1093/qmath/haad031}.
\end{thebibliography}
\end{document}
